In school $s$, potential class size varies depending on whether the school is assigned to the tracking condition or not. In a school $s$, let $N_{zsc}$ represent the class size in school $s$ and class $c$ under tracking condition $z$.  
\[N_{1sc} = \sum_{i \in \mathcal{I}_s} \mathbbm{1}\{ r_i \in \bar{r}_c \} \]
\[N_{0sc} = \sum_{i \in \mathcal{I}_s} \mathbbm{1}\{ a_i=c \} \]
\[N_{1sci} = \sum_{c \in \{0,1\}} N_{1sc} \mathbbm{1}\{r_i \in \bar{r}_c\} \]
\[N_{0sci} = \sum_{c \in \{0,1\}} N_{0sc} \mathbbm{1}\{a_i = c\} \]
where $r_i$ is the assignment score, and $a_i$ is the baseline grade of student $i$. Given a probability $p_{sci} = E[Z_{sci}|s]$, the expected classize for individual $i$ is therefore
\[ \mu_{Ni} = p_{sci} N_{0sci} + (1-p_{sci}) N_{1sci} \]
and the actual class size for student $i$ reflects the realized tracking status
\[ N_{sci} = N_{0sci} + Z_i ( N_{1sci} - N_{0sci} )\]

We want to estimate the following linear structural equation for test scores as a function of class size
\[ Y_{sci} = \beta N_{sci} + X_{sci}' \Gamma + \varepsilon_i \]
We are able to estimate 
\[ Y_{sci} = \beta N_{sci} + \lambda \mu_{Ni} + \eta_i \]
Identification is a consequence of the following conditional independence property:
\[E[Y_{0i}|\mu_{Ni},N_{sci}] = E\left[Y_{0i} | Z_i\right] \]
\[ W_{sci} = N_{0sc} + Z_i ( )\]

To measure the overall impact of the intervention on test scores, we will run regressions of the form
\[ y_{ij} = \alpha T_j + X_{ij} \beta + \epsilon_{ij} \]
Where $y_{ij}$ is a test score outcome $T_j \in \{0,1\}$ indicates  treatment assignment of school $j$, and $X_{ij}$ is a vector of individual and school level characteristics, as well as a constant. $X_{ij}$ may include baseline test score controls.

Because the intervention was conducted for both Standards 1 and 2, we will also explore whether the intervention's impact was different in each grade.

To identify the potential differential effects for children assigned to the lower (blue) and upper (purple) group, we also will run specifications of the form

\[ y_{ij} = \alpha T_j + \gamma T_j \times P_{i} + X_{ij} \beta + \epsilon_{ij} \]
Where $P_i$ indicates that a student passed the grade specific threshold for entry into the upper (purple) group. 